\chapter{Casos de uso}

El diagrama de casos de uso (\cref{fig:casos-uso}) muestra qué puede hacer cada actor. El Superintendente registra el catálogo de conceptos y las estimaciones. El Supervisor elabora los reportes diarios y mensuales y revisa las estimaciones del Superintendente. El Residente valida los reportes y da la validación final de las estimaciones. El Administrador crea las cuentas y puede cambiar el estatus de las validaciones. El Directivo solo consulta. Las notas de bitácora, la carga de documentos y la consulta del proyecto las comparten varios actores.

\begin{figure}[h]
\centering
\resizebox{\linewidth}{!}{%
\begin{tikzpicture}[
  cu/.style={ellipse,draw=guinda!70,fill=guindaclaro,line width=0.5pt,text width=2.9cm,align=center,
             font=\sffamily\scriptsize,inner sep=1.2pt,minimum height=0.95cm},
  actor/.style={font=\sffamily\scriptsize\bfseries,align=center},
  enlace/.style={draw=black!45,line width=0.35pt}]
  % actores (monigote)
  \newcommand{\monigote}[3]{%
    \begin{scope}[shift={(#1)}]
      \draw[line width=0.7pt] (0,0.55) circle (0.18);
      \draw[line width=0.7pt] (0,0.37) -- (0,-0.15) (-0.3,0.2) -- (0.3,0.2) (0,-0.15) -- (-0.22,-0.55) (0,-0.15) -- (0.22,-0.55);
      \node[actor,below] at (0,-0.6) {#2};
      \coordinate (#3) at (0,0);
    \end{scope}}
  \monigote{-8.2,-2.4}{Superintendente\\(contratista)}{SI}
  \monigote{-8.2,-8.4}{Residente}{RE}
  \monigote{-8.2,-13.2}{Directivo}{DI}
  \monigote{8.2,-2.4}{Supervisor}{SU}
  \monigote{8.2,-12.2}{Administrador}{AD}
  % límite del sistema
  \begin{pgfonlayer}{background}
    \draw[rounded corners=6pt,draw=black!35,fill=grisfondo] (-5.6,0.3) rectangle (5.6,-15.6);
  \end{pgfonlayer}
  \node[font=\sffamily\small\bfseries,anchor=north west,text=black!60] at (-5.45,0.2) {Sistema de Obras};
  % columna izquierda
  \node[cu] (c01) at (-3.6,-1.1) {CU-01 Crear catálogo de conceptos};
  \node[cu] (c02) at (-3.6,-2.5) {CU-02 Crear estimaciones};
  \node[cu] (c06) at (-3.6,-7.3) {CU-06 Validar reportes diarios};
  \node[cu] (c07) at (-3.6,-8.7) {CU-07 Validar reportes mensuales};
  \node[cu] (c08) at (-3.6,-10.1) {CU-08 Validar estimaciones del Supervisor};
  % columna central
  \node[cu] (c09) at (0,-3.9) {CU-09 Hacer notas de bitácora};
  \node[cu] (c10) at (0,-5.3) {CU-10 Firmar notas de bitácora};
  \node[cu] (c11) at (0,-6.7) {CU-11 Subir documentos al repositorio};
  \node[cu] (c12) at (0,-11.9) {CU-12 Consultar información del proyecto};
  % columna derecha
  \node[cu] (c03) at (3.6,-1.1) {CU-03 Crear reportes diarios};
  \node[cu] (c04) at (3.6,-2.5) {CU-04 Crear reportes mensuales};
  \node[cu] (c05) at (3.6,-3.9) {CU-05 Validar estimaciones del Superintendente};
  \node[cu] (c13) at (3.6,-8.7) {CU-13 Crear cuentas para cada rol};
  \node[cu] (c14) at (3.6,-10.1) {CU-14 Cambiar estatus de reportes diarios};
  \node[cu] (c15) at (3.6,-11.5) {CU-15 Cambiar estatus de reportes mensuales};
  \node[cu] (c16) at (3.6,-12.9) {CU-16 Cambiar estatus de estimaciones (Supervisor)};
  \node[cu] (c17) at (3.6,-14.3) {CU-17 Cambiar estatus de estimaciones (Residente)};
  % enlaces: entre el recuadro (fondo) y los casos de uso
  \begin{pgfonlayer}{lineas}
  \foreach \c in {c01,c02,c09,c10,c11,c12} \draw[enlace] ([xshift=0.35cm]SI) -- (\c);
  \foreach \c in {c06,c07,c08,c09,c10,c11,c12} \draw[enlace] ([xshift=0.35cm]RE) -- (\c);
  \draw[enlace] ([xshift=0.35cm]DI) -- (c12);
  \foreach \c in {c03,c04,c05,c09,c10,c11,c12} \draw[enlace] ([xshift=-0.35cm]SU) -- (\c);
  \foreach \c in {c12,c13,c14,c15,c16,c17} \draw[enlace] ([xshift=-0.35cm]AD) -- (\c);
  \end{pgfonlayer}
\end{tikzpicture}}
\caption{Diagrama de casos de uso}\label{fig:casos-uso}
\end{figure}


\section{Índice de casos de uso}
Del diagrama se obtienen 17 casos de uso (\cref{tab:indice-cu}). Los identificadores siguen el orden de las funciones de cada rol, del Superintendente al Administrador; las funciones compartidas se describen una sola vez.

\begin{table}[h]
\centering\small
\caption{Índice de casos de uso}\label{tab:indice-cu}
\renewcommand{\arraystretch}{1.12}
\rowcolors{2}{grisfondo}{white}
\begin{tabularx}{\linewidth}{@{}>{\sffamily\bfseries\color{guinda}}lX>{\raggedright\arraybackslash}p{5.4cm}@{}}
\toprule
\normalfont\sffamily\bfseries ID & \sffamily\bfseries Caso de uso & \sffamily\bfseries Actores\\\midrule
CU-01 & Crear catálogo de conceptos & Superintendente\\
CU-02 & Crear estimaciones & Superintendente\\
CU-03 & Crear reportes diarios & Supervisor\\
CU-04 & Crear reportes mensuales & Supervisor\\
CU-05 & Validar estimaciones del Superintendente & Supervisor\\
CU-06 & Validar reportes diarios & Residente\\
CU-07 & Validar reportes mensuales & Residente\\
CU-08 & Validar estimaciones del Supervisor & Residente\\
CU-09 & Hacer notas de bitácora & Superintendente, Supervisor, Residente\\
CU-10 & Firmar notas de bitácora & Superintendente, Supervisor, Residente\\
CU-11 & Subir documentos al repositorio & Superintendente, Supervisor, Residente\\
CU-12 & Consultar información del proyecto & Todos los roles\\
CU-13 & Crear cuentas para cada rol & Administrador\\
CU-14 & Cambiar estatus de validaciones de reportes diarios & Administrador\\
CU-15 & Cambiar estatus de validaciones de reportes mensuales & Administrador\\
CU-16 & Cambiar estatus de validaciones de estimaciones del Superintendente & Administrador\\
CU-17 & Cambiar estatus de validaciones de estimaciones del Residente & Administrador\\
\bottomrule
\end{tabularx}
\end{table}

\section{Descripción de los casos de uso}
Cada caso de uso se describe con el mismo formato: actores, requerimientos que atiende, descripción, flujo principal (el recorrido normal), flujos alternativos (datos incompletos, rechazos, atrasos o cambios autorizados), precondiciones y postcondiciones.

\casouso{CU-01}{Superintendente}{Crear catálogo de conceptos}
{Permite registrar el catálogo de conceptos de la obra (clave, capítulo, descripción, unidad, cantidad y precio unitario), que sirve de base para el seguimiento de los trabajos y las estimaciones.}
{\item El Superintendente abre el catálogo de la obra.
 \item El sistema muestra el contrato y la tabla de captura.
 \item El Superintendente captura cada concepto; el sistema calcula el importe de cada renglón y el total.
 \item El Superintendente confirma el registro.
 \item El sistema valida los datos, guarda el catálogo y lo asocia a la obra; muestra la diferencia entre el total del catálogo y el monto del contrato.}
{\item \textbf{4a. Datos incompletos o inválidos:} el sistema marca los campos con error y no guarda hasta que se corrijan.
 \item \textbf{4b. Clave repetida:} el sistema avisa que la clave ya existe y solicita otra.
 \item \textbf{Catálogo en uso:} si ya hay estimaciones con el catálogo, no se puede modificar; los cambios se tramitan como reprogramación o trabajo adicional, sin modificar los montos del contrato.}
{\item Sesión activa con rol de Superintendente.
 \item La obra y su contrato están registrados.}
{\item El catálogo queda guardado y asociado a la obra.
 \item Los conceptos quedan disponibles para las estimaciones.}
{RF-07, RF-08}

\casouso{CU-02}{Superintendente}{Crear estimaciones}
{Permite registrar una estimación de trabajos ejecutados en un periodo, con sus documentos de soporte, para que la revise el Supervisor y la valide el Residente.}
{\item El Superintendente elige crear una estimación; el sistema propone el periodo siguiente a la última estimación.
 \item El sistema muestra los conceptos del catálogo con lo contratado y lo estimado antes.
 \item El Superintendente captura las cantidades ejecutadas; el sistema calcula en pantalla importes y acumulados.
 \item Al guardar, el sistema calcula importe, amortización del anticipo, IVA, retención por derechos y, si la obra va atrasada a la fecha de corte, la retención por atraso (\cref{cap:reglas}). La estimación queda como \estatus{Borrador} con folio \cod{EST-001}, \cod{EST-002}\ldots
 \item El Superintendente adjunta los documentos de soporte; los números generadores son obligatorios.
 \item El Superintendente revisa los importes y la envía; queda \estatus{Pendiente} y se avisa al Supervisor.}
{\item \textbf{3a. Concepto excedente:} si el acumulado rebasa lo contratado, el renglón debe citar la nota de bitácora firmada que autorizó el excedente (\reg, art.~118).
 \item \textbf{3b. Atraso no imputable al contratista:} el Superintendente lo indica con su justificación y evidencia; no se aplica la retención por atraso y queda a revisión de los validadores.
 \item \textbf{Periodo mayor a un mes o traslapado:} el sistema no lo acepta (\ley, art.~54).
 \item \textbf{Falta soporte:} sin números generadores el sistema no permite enviarla.
 \item \textbf{Estimación rechazada o con observaciones:} el Superintendente la corrige y la reenvía; las decisiones previas se reinician.}
{\item Sesión activa con rol de Superintendente.
 \item El contrato está formalizado y tiene catálogo de conceptos.
 \item No hay otra estimación abierta (sin aprobar) del contrato.}
{\item La estimación queda \estatus{Pendiente} y se avisa al Supervisor.
 \item Sus volúmenes quedan bloqueados mientras está en validación.}
{RF-10, RF-14}

\casouso{CU-03}{Supervisor}{Crear reportes diarios}
{Permite registrar el reporte diario de la obra con el avance físico real, la jornada, el clima, las actividades, las incidencias, las notas de bitácora del día y las evidencias.}
{\item El Supervisor abre los reportes diarios de la obra y elige uno nuevo.
 \item El sistema muestra el formulario con la fecha del día.
 \item El Supervisor captura la jornada (inicio y fin), el clima, la temperatura, el avance físico real, las actividades y las observaciones, y elige en un menú desplegable las notas de bitácora relacionadas.
 \item El sistema calcula la desviación contra el programa (atraso o adelanto, en puntos y en días) y el resumen financiero.
 \item El Supervisor guarda el borrador, adjunta evidencias y lo envía a revisión.
 \item El sistema lo guarda como \estatus{Pendiente}, avisa al Residente y permite descargar el PDF.}
{\item \textbf{4a. Atraso crítico:} si el avance va un día o más detrás del programa, el sistema exige documentar la causa y si se solicita reprogramación o retención.
 \item \textbf{6a. Información incompleta o incorrecta:} el sistema indica cada campo con error (jornada, clima, temperatura entre $-20$ y 55\,°C, actividades, avance entre 0 y 100\,\%) y no guarda hasta que se corrija.
 \item \textbf{Fecha futura o repetida:} en cuanto se elige la fecha, el sistema avisa si es posterior al día actual, anterior al inicio de la obra o si ya existe un reporte de la obra para ese día, y no lo guarda.
 \item \textbf{Avance menor al anterior:} el avance físico es acumulado; no puede ser menor que el de un reporte de un día anterior ni mayor que el de uno posterior.}
{\item Sesión activa con rol de Supervisor (el Administrador no elabora reportes).
 \item La obra está registrada; para calcular la desviación, tiene programa de obra.}
{\item El reporte queda en el historial de la obra y se avisa a los validadores.
 \item Al aprobarse, el avance físico actualiza el seguimiento del contrato.}
{RF-09, RF-11, RF-16, RF-22}

\casouso{CU-04}{Supervisor}{Crear reportes mensuales}
{Permite consolidar los reportes diarios de un mes en un reporte mensual con el avance acumulado, las estimaciones y la bitácora del periodo.}
{\item El Supervisor elige el mes a reportar.
 \item El sistema consolida los reportes diarios, el avance físico y financiero acumulado, las estimaciones y las notas de bitácora del periodo, y muestra las advertencias.
 \item El Supervisor agrega el resumen, las observaciones y las evidencias.
 \item El Supervisor revisa la vista previa y lo envía.
 \item El sistema lo guarda como \estatus{Pendiente}, avisa a los validadores y permite descargar el PDF. Al aprobarse, las cifras quedan congeladas.}
{\item \textbf{2a. Faltan reportes diarios en el periodo:} el sistema lo advierte con los días faltantes; el Supervisor decide si los completa antes de continuar.
 \item \textbf{3a. Información incompleta:} el Supervisor completa los datos o guarda el reporte como borrador.}
{\item Sesión activa con rol de Supervisor.
 \item Existen reportes diarios registrados en el mes.}
{\item El reporte mensual queda \estatus{Pendiente} y se avisa a los validadores.}
{RF-16, RF-19}

\casouso{CU-05}{Supervisor}{Validar estimaciones del Superintendente}
{Permite al Supervisor revisar las estimaciones enviadas por el Superintendente y aprobarlas, regresarlas con observaciones o rechazarlas, antes de que pasen al Residente.}
{\item El Supervisor abre una estimación \estatus{Pendiente} desde su bandeja.
 \item El sistema muestra conceptos, cantidades, importes, documentos de soporte, avance y plazos de revisión.
 \item El Supervisor verifica la información contra el avance real y la bitácora.
 \item El Supervisor elige el resultado y escribe sus comentarios.
 \item El sistema registra la validación con usuario, rol, fecha y comentarios.}
{\item \textbf{4a. Con observaciones:} la estimación regresa al Superintendente para corregirla sin que cuente como rechazo.
 \item \textbf{4b. Rechazada:} regresa al Superintendente con el motivo.
 \item \textbf{Plazo vencido:} el sistema avisa que pasaron los 15 días de revisión, sin autorizarla en automático.}
{\item Sesión activa con rol de Supervisor.
 \item Existe una estimación \estatus{Pendiente}.}
{\item Aprobada, queda \estatus{En revisión} y se avisa al Superintendente y al Residente.
 \item Con observaciones o rechazada, se avisa al Superintendente.}
{RF-15}

\casouso{CU-06}{Residente}{Validar reportes diarios}
{Permite al Residente revisar los reportes diarios del Supervisor y aprobarlos o rechazarlos; decide una sola vez.}
{\item El validador abre un reporte \estatus{Pendiente} desde su bandeja.
 \item El sistema muestra avance, desviación, actividades, notas vinculadas y evidencias.
 \item El validador revisa la información.
 \item El validador aprueba o rechaza con comentarios.
 \item El sistema registra la decisión con usuario, rol, fecha y comentarios.}
{\item \textbf{4a. Reporte rechazado:} regresa al Supervisor con el motivo; al corregirlo vuelve a revisión.}
{\item Sesión activa con rol de Residente.
 \item Existe un reporte diario enviado a revisión.}
{\item El reporte queda \estatus{Aprobado} o \estatus{Rechazado}.
 \item Se avisa al Supervisor del resultado.}
{RF-17}

\casouso{CU-07}{Residente}{Validar reportes mensuales}
{Permite al Residente revisar y validar los reportes mensuales del Supervisor.}
{\item El validador abre un reporte mensual pendiente.
 \item El sistema muestra el consolidado del mes, las observaciones y las evidencias.
 \item El validador revisa la información.
 \item El validador aprueba o rechaza con comentarios.
 \item El sistema registra la decisión con usuario, rol, fecha y comentarios.}
{\item \textbf{4a. Reporte rechazado:} regresa al Supervisor con el motivo para su corrección.}
{\item Sesión activa con rol de Residente.
 \item Existe un reporte mensual enviado a revisión.}
{\item El reporte queda \estatus{Aprobado} o \estatus{Rechazado}; al aprobarse se congelan sus cifras.
 \item Se avisa al Supervisor y, si se aprueba, a los Directivos.}
{RF-17}

\casouso{CU-08}{Residente}{Validar estimaciones del Supervisor}
{Permite al Residente dar la validación final a las estimaciones ya aprobadas por el Supervisor.}
{\item El Residente abre una estimación \estatus{En revisión} desde su bandeja.
 \item El sistema muestra la estimación, la revisión del Supervisor y sus comentarios.
 \item El Residente revisa la información.
 \item El Residente aprueba o rechaza con comentarios.
 \item El sistema registra la validación con usuario, rol, fecha y comentarios.}
{\item \textbf{4a. Estimación rechazada:} regresa al Superintendente con el motivo.}
{\item Sesión activa con rol de Residente.
 \item La estimación fue aprobada previamente por el Supervisor.}
{\item La estimación queda \estatus{Aprobado} o \estatus{Rechazado}; aprobada, cuenta para el avance financiero.
 \item Se avisa al Superintendente y al Supervisor.}
{RF-15}

\casouso{CU-09}{Superintendente, Supervisor, Residente}{Hacer notas de bitácora}
{Permite registrar notas en la bitácora de la obra (instrucciones, solicitudes, avances, incidencias y otras), con su tipo, nota relacionada y evidencias.}
{\item El usuario abre la bitácora de la obra y elige una nota nueva.
 \item El usuario captura el tipo, la descripción y, si responde a otra, la nota relacionada.
 \item El usuario confirma el registro.
 \item El sistema asigna el folio consecutivo, guarda la nota firmada por su autor y queda pendiente de las demás firmas.}
{\item \textbf{3a. Datos incompletos:} el sistema indica la información faltante y no guarda la nota.
 \item \textbf{3b. Evidencia no permitida:} el sistema rechaza el archivo e indica el motivo.
 \item \textbf{Corrección:} el autor puede editarla o borrarla mientras nadie más la haya firmado; el folio borrado no se reutiliza.}
{\item Sesión activa con rol de Superintendente, Supervisor o Residente.
 \item La obra está registrada.}
{\item La nota queda en la bitácora con estado \estatus{Pendiente}.
 \item Se avisa a los demás firmantes.}
{RF-12, RF-22}

\casouso{CU-10}{Superintendente, Supervisor, Residente}{Firmar notas de bitácora}
{Permite firmar o rechazar una nota de bitácora para dejar constancia de que se revisó y si se está de acuerdo con su contenido.}
{\item El usuario consulta las notas pendientes de firma.
 \item El usuario abre una nota y revisa su contenido.
 \item El usuario firma o la rechaza con motivo.
 \item El sistema registra la firma con usuario, rol, fecha y hora, y actualiza el estado.}
{\item \textbf{3a. Usuario no autorizado:} el sistema no permite firmar e informa el motivo.
 \item \textbf{3b. Nota rechazada:} regresa al autor; al corregirla se reinician las firmas.
 \item \textbf{Firma pendiente:} la nota permanece \estatus{Pendiente} y sigue en la bandeja de quien debe firmar.}
{\item Sesión activa con rol de Superintendente, Supervisor o Residente.
 \item Existe una nota pendiente de firma.}
{\item Con todas las firmas queda \estatus{Firmada} y ya no puede modificarse.
 \item Cada firma queda en el historial.}
{RF-13}

\casouso{CU-11}{Superintendente, Supervisor, Residente}{Subir documentos al repositorio}
{Permite cargar documentos y archivos (contratos, reportes, documentos técnicos, planos, fotografías, facturas y evidencias) y asociarlos a la obra.}
{\item El usuario abre el repositorio de documentos de la obra.
 \item El usuario elige el archivo y captura tipo, nombre y descripción.
 \item El usuario confirma la carga.
 \item El sistema valida el formato y el tamaño, guarda el archivo y lo asocia a la obra con usuario y fecha.}
{\item \textbf{3a. Archivo no permitido o demasiado grande:} el sistema rechaza la carga e indica el motivo (se aceptan PDF, Word, Excel, imágenes y DWG de hasta 20\,MB, reconocidos por su contenido).
 \item \textbf{3b. Falla de carga:} el sistema informa el error y permite reintentar.
 \item \textbf{Borrar un documento:} solo quien lo subió o el Administrador; queda registrado en el historial.}
{\item Sesión activa con rol de Superintendente, Supervisor o Residente.
 \item La obra está registrada.}
{\item El documento queda disponible en el repositorio de la obra.
 \item La carga queda registrada en el historial de la obra.}
{RF-06, RF-20, RF-25}

\casouso{CU-12}{Todos los roles}{Consultar información del proyecto}
{Permite consultar el expediente de una obra: contrato, catálogo, reportes, estimaciones, bitácora, avances y documentos. Es el único caso de uso del Directivo, que solo consulta.}
{\item El usuario ingresa al módulo de obras y busca por nombre, identificador, contrato, contratista o estado.
 \item El sistema muestra las obras que coinciden con su avance real contra el programado.
 \item El usuario selecciona una obra.
 \item El sistema muestra el expediente con los indicadores y la gráfica de avance físico y financiero (programado contra real), según los permisos del rol.}
{\item \textbf{2a. Sin resultados:} el sistema informa que no hay coincidencias.
 \item \textbf{4a. Información restringida para el rol:} el sistema oculta lo que el usuario no puede ver (por ejemplo, borradores ajenos).
 \item \textbf{4b. Error al cargar la información:} la sección afectada muestra un aviso y permite reintentar sin afectar al resto.}
{\item Sesión activa.
 \item Existen obras con información capturada.}
{\item El usuario visualiza la información; no se modifican datos.}
{RF-02, RF-05, RF-11, RF-18, RF-21}

\casouso{CU-13}{Administrador}{Crear cuentas para cada rol}
{Permite crear cuentas, asignar su rol en cada obra y definir la contraseña inicial; también modificar o desactivar cuentas.}
{\item El Administrador abre el módulo de usuarios y elige una cuenta nueva.
 \item Captura nombre, apellidos, correo, contraseña inicial y, opcionalmente, la obra y el rol.
 \item Confirma.
 \item El sistema valida los datos, guarda la cuenta, aplica los permisos del rol y registra la acción.}
{\item \textbf{3a. Correo ya registrado o datos inválidos:} el sistema indica el error y no guarda la cuenta.
 \item \textbf{Modificar o desactivar:} el Administrador cambia datos, roles o estado; una cuenta desactivada ya no puede iniciar sesión.}
{\item Sesión activa con rol de Administrador.}
{\item La cuenta queda creada o actualizada con los permisos de su rol.
 \item La acción queda en el historial de la cuenta.}
{RF-23, RF-24}

\casouso{CU-14}{Administrador}{Cambiar estatus de validaciones de reportes diarios}
{Permite cambiar el estatus de la validación de un reporte diario (Pendiente, Aprobado o Rechazado, por rol) sin modificar su contenido.}
{\item El Administrador abre el módulo de validaciones, pestaña de reportes diarios.
 \item El sistema muestra los reportes con la decisión vigente de cada validador.
 \item El Administrador elige un reporte, el rol y el nuevo estatus.
 \item Captura el motivo y confirma.
 \item El sistema registra la decisión con usuario, fecha y motivo, y recalcula el estatus del reporte.}
{\item \textbf{4a. Falta el motivo:} el sistema no permite continuar hasta capturarlo.}
{\item Sesión activa con rol de Administrador.
 \item Existe al menos un reporte enviado a revisión.}
{\item El estatus queda actualizado y el cambio en el historial.
 \item Se avisa al Supervisor y a los validadores.}
{RF-25}

\casouso{CU-15}{Administrador}{Cambiar estatus de validaciones de reportes mensuales}
{Permite cambiar el estatus de la validación de un reporte mensual sin modificar su contenido; al regresarlo a revisión se liberan sus cifras congeladas.}
{\item El Administrador abre el módulo de validaciones, pestaña de reportes mensuales.
 \item El sistema muestra los reportes con la decisión vigente de cada validador.
 \item El Administrador elige un reporte, el rol y el nuevo estatus.
 \item Captura el motivo y confirma.
 \item El sistema registra la decisión con usuario, fecha y motivo, y recalcula el estatus.}
{\item \textbf{4a. Falta el motivo:} el sistema no permite continuar hasta capturarlo.}
{\item Sesión activa con rol de Administrador.
 \item Existe al menos un reporte mensual enviado a revisión.}
{\item El estatus queda actualizado y el cambio en el historial.
 \item Se avisa al Supervisor y a los validadores.}
{RF-25}

\casouso{CU-16}{Administrador}{Cambiar estatus de validaciones de estimaciones del Superintendente}
{Permite cambiar la decisión que el Supervisor registró sobre una estimación (Pendiente, Aprobado, Con observaciones o Rechazado), sin modificar la estimación.}
{\item El Administrador abre el módulo de validaciones, pestaña de estimaciones.
 \item El sistema muestra las estimaciones con la decisión del Supervisor y del Residente.
 \item El Administrador elige una estimación, el rol Supervisor y el nuevo estatus.
 \item Captura el motivo y confirma.
 \item El sistema registra la decisión con usuario, fecha y motivo, y recalcula el estatus.}
{\item \textbf{4a. Falta el motivo:} el sistema no permite continuar hasta capturarlo.}
{\item Sesión activa con rol de Administrador.
 \item Existe al menos una estimación enviada a revisión.}
{\item El estatus queda actualizado y el cambio en el historial.
 \item Se avisa al Superintendente, al Supervisor y al Residente.}
{RF-25}

\casouso{CU-17}{Administrador}{Cambiar estatus de validaciones de estimaciones del Residente}
{Permite cambiar la decisión que el Residente registró sobre una estimación, sin modificar la estimación.}
{\item El Administrador abre el módulo de validaciones, pestaña de estimaciones.
 \item El sistema muestra las estimaciones con la decisión de cada validador.
 \item El Administrador elige una estimación, el rol Residente y el nuevo estatus.
 \item Captura el motivo y confirma.
 \item El sistema registra la decisión con usuario, fecha y motivo, y recalcula el estatus.}
{\item \textbf{4a. Falta el motivo:} el sistema no permite continuar hasta capturarlo.}
{\item Sesión activa con rol de Administrador.
 \item Existe al menos una estimación con decisión del Residente.}
{\item El estatus queda actualizado y el cambio en el historial.
 \item Se avisa al Residente, al Superintendente y al Supervisor.}
{RF-25}
